\section{The PyMORGAN boundary}\label{sec:boundary}

Most questions of the form ``where does this code belong?'' are answered by
Table~\ref{tab:boundary}. Read it before adding a module.

\begin{table}[h]
  \centering
  \begin{tabular}{lll}
    \toprule
    Concern & Owner & Entry point \\
    \midrule
    File formats, loader registry      & PyMORGAN & \code{pm.load\_1D}, \code{pm.register\_loader} \\
    Dataset objects                    & PyMORGAN & \code{pm.Dataset1D}, \code{pm.Dataset2D} \\
    Background, chirp, solvent         & PyMORGAN & \code{Dataset1D.background\_correct}, \dots \\
    Contour / spectra / kinetics plots & PyMORGAN & \code{Dataset1D.plot\_*} \\
    Species-spectra rendering          & PyMORGAN & \code{plot\_species\_spectra} \\
    Aesthetics, colourmaps, unit labels & PyMORGAN & \code{pm.Settings}, \code{pm.apply\_style()} \\
    \midrule
    Kinetic models, rate matrices      & PyRATE-TA & \code{pyrate_ta.models} \\
    Fitting engines, global \& target  & PyRATE-TA & \code{pyrate_ta.fit} \\
    Fit results, uncertainties         & PyRATE-TA & \code{pyrate_ta.results} \\
    Analysis-only plots                & PyRATE-TA & \code{pyrate_ta.plot} \\
    \bottomrule
  \end{tabular}
  \caption{Ownership of each concern across the two packages.}
  \label{tab:boundary}
\end{table}

Three rules follow.

\paragraph{The dependency is one-way.}
\code{import pymorgan} inside PyRATE-TA is expected; \code{import pyrate_ta} inside
PyMORGAN is a design error. PyMORGAN's \code{plot\_species\_spectra}
deliberately takes plain arrays (\code{Sfit}, \code{Taus}, \code{TauErr},
\code{isFixTau}, \code{modelType}) rather than a PyRATE-TA result object, precisely
so the arrow never has to reverse; the adapter that produces those arrays lives
on the PyRATE-TA side. \code{scripts/run\_checks.py} asserts the direction.

\paragraph{No file readers here.}
Loading goes through \code{pm.load\_1D(path, data\_type=\dots)}, which resolves
the format through PyMORGAN's registry and therefore picks up every format
PyMORGAN supports, present and future. A new format is contributed upstream to
\code{pymorgan.oneD.load}, or registered at runtime with
\code{pm.register\_loader} --- never forked into PyRATE-TA.

\paragraph{One aesthetics system.}
Call \code{pm.apply\_style()} and read \code{pm.get\_settings()} for label
style, colourmap, time-axis convention and figure sizes. PyRATE-TA's own
\code{Settings} holds analysis-only fields (chapter~\ref{sec:settings}).

\subsection{A minimal session}

\begin{lstlisting}
import pymorgan as pm
import pyrate_ta as pr

data = pm.load_1D("scan.pdat", data_type="PDAT")  # PyMORGAN loads
data.background_correct(tmin=-20, tmax=-5)        # PyMORGAN processes
fit  = pr.fit_global(data, n_components=3)        # PyRATE-TA fits
data.plot_species_spectra(*fit.as_species_args()) # PyMORGAN plots
\end{lstlisting}

\subsection{Units}\label{sec:units}

Signal data is in mOD, delays are in the dataset's own time unit and the probe
axis in its native spectral unit. Fitted lifetimes are reported in the
dataset's time unit and are never silently converted: a lifetime printed
without its unit, or converted behind the user's back, is a bug.

\subsection{Extraction of traces}

PyMORGAN removed its own \code{extract\_kinetic}/\code{fit\_kinetics} when
PyRATE-TA took over analysis, so there is no upstream extraction helper to call:
that primitive lives here. For a headless fit, slice the dataset directly
(\code{data.Z[:, idx, detector]}) rather than re-deriving pixel-lookup logic;
\code{Dataset1D.plot\_kinetics} returns the extracted data as its third value
when a figure is wanted as well.
